\begingroup
\widowpenalty=10000
\SourceLeaf{112}
adeoque (prop. 10 art. 132) $+fRp$. Sed est etiam $+faRp$, quare fiet $+aRp$, sive $-aNp$.

II. Quando $e$ per $p$ est divisibilis, sit $e=pg$, atque $f=ph$. Erit itaque $g^2p=1+ha$. Tum $h$ erit positivus, formae $4n+3$ $(B)$, et ad $p$ et $g^2$ primus. Porro $+g^2pRh$, adeoque $+pRh$; hinc fit (prop. 13 art. 132) $-hRp$. At est $-haRp$, unde fit $+aRp$ atque $-aNp$.

\Article{139}

\textit{Casus tertius. Quando $T+1$ est formae $4n+1$ $(=a)$, $p$ eiusdem formae, atque $\pm pNa$: non potest esse $\pm aRp$.} (Supra casus secundus).

Capiatur aliquis numerus primus ipso $a$ minor, cuius non-residuum sit $+a$, quales dari supra demonstravimus (artt. 125, 129). Sed hic duos casus seorsim considerare oportet, prout hic numerus primus fuerit formae $4n+1$ vel $4n+3$, non enim demonstratum fuit, dari tales numeros primos \textit{utriusque} formae.

I. Sit iste numerus primus formae $4n+1$ et $=a'$. Tum erit $\pm a'Na$ (art. 131) adeoque $\pm a'pRa$. Sit igitur $e^2\equiv a'p\ (\mathrm{mod.}\ a)$ atque $e$ par, $<a$. Tunc iterum quatuor casus erunt distinguendi.

1) Quando $e$ neque per $p$ neque per $a'$ est divisibilis. Ponatur $e^2=a'p\pm af$, signis ita acceptis ut $f$ fiat positivus. Tum erit $f<a$, ad $a'$ et $p$ primus atque pro signo superiori formae $4n+3$, pro inferiori formae $4n+1$. Designemus brevitatis gratia per $[x,y]$ multitudinem factorum primorum numeri $y$ quorum non-residuum est $x$. Tum erit $a'pRf$ adeoque $[a'p,f]=0$. Hinc erit $[f,a'p]$ numerus par (propp. 1, 3, art. 133), \textit{i. e.} aut $=0$ aut $=2$. Quare erit $f$ aut residuum utriusque numerorum $a',p$, aut neutrius. Illud autem est impossibile, quum $\pm af$ sit residuum ipsius $a'$, atque $\pm aNa'$ (hyp.); unde fit $\pm fNa'$. Hinc $f$ debet esse utriusque numerorum $a',p$ non-residuum. At propter $\pm afRp$ erit $\pm aNp$. \textit{Q.~E.~D.}

2) Quando $e$ per $p$, neque vero per $a'$ est divisibilis, sit $e=gp$, atque $g^2p=a'\pm ah$, signo ita determinato, ut $h$ fiat positivus. Tum erit $h<a$, ad $a',g$ et $p$ primus. atque pro signo superiori formae $4n+3$, pro inferiori vero formae $4n+1$. Ex aequatione $g^2p=a'\pm ah$ si per $p$ et $a'$ multiplicatur, nullo negotio deduci potest, $pa'Rh\,.\,.\,.\,.\,.(\alpha)$; $\pm ahpRa'\,.\,.\,.\,.(\beta)$; $aa'hRp\,.\,.\,.\,.\,.(\gamma)$. Ex $(\alpha)$ sequitur $[pa',h]=0$, adeoque (propp. 1, 3, art. 133) $[h,pa']$ par, \textit{i. e.}
\SourceLeaf{113}
erit $h$ non-residuum vel utriusque $p,a'$, vel neutrius. \textit{Priori in casu} ex $(\beta)$ sequitur $\pm apNa'$, et quum per hyp. sit $\pm aNa'$, erit $\pm pRa'$. Hinc per theor. fundam. quod pro numeris $p,a'$ ipso $T+1$ minoribus valet, $\pm a'Rp$. Hinc et ex eo quod $hNp$, fit per $(\gamma)$ $\pm aNp$. \textit{Q.~E.~D.} \textit{Posteriori casu} ex $(\beta)$ sequitur $\pm apRa'$, hinc $\pm pNa'$, $\pm a'Np$, hincque tandem et ex $hRp$ fit ex $(\gamma)$ $\pm aNp$. \textit{Q.~E.~D.}

3) Quando $e$ per $a'$ non autem per $p$ est divisibilis. Pro hoc casu demonstratio tantum non eodem modo procedit ut in praec., neminemque qui hanc penetravit poterit morari.

4) Quando $e$ tum per $a'$ tum per $p$ est divisibilis adeoque etiam per productum $a'p$ (numeros $a',p$ enim \textit{inaequales} esse supponimus, quia alias id quod demonstrare operam damus, esse $aNp$ iam in hypothesi $aNa'$ contentum foret), sit $e=ga'p$ atque $g^2a'p=1\pm ah$. Tum erit $h<a$, ad $a'$ et $p$ primus atque pro signo superiori formae $4n+3$, pro inferiori formae $4n+1$. Facile vero perspicitur, ex ista aequatione deduci posse haec $a'pRh\,.\,.\,.\,.(\alpha)$; $\pm ahRa'\,.\,.\,.\,.(\beta)$; $\pm ahRp\,.\,.\,.\,.(\gamma)$. Ex $(\alpha)$ quod convenit cum $(\alpha)$ in (2) sequitur perinde ut illic, esse vel simul $hRp$, $hRa'$, vel $hNp$, $hNa'$. Sed in casu priori foret per $(\beta)$, $aRa'$, contra hyp.; quare erit $hNp$, adeoque per $(\gamma)$ etiam $aNp$.

II. Quando iste numerus primus est formae $4n+3$, demonstratio praecedenti tam similis est, ut eam apponere superfluum nobis visum sit. In eorum gratiam qui per se eam evolvere gestiunt (quod maxime commendamus), id tantum observamus, postquam ad talem aequationem $e^2=bp\pm af$ (designante $b$ illum numerum primum) perventum fuerit, ad perspicuitatem profuturum, si utrumque signum seorsim consideretur.

\Article{140}

\textit{Casus quartus. Quando $T+1$ est formae $4n+1$ $(=a)$, $p$ formae $4n+3$, atque $\pm pNa$, non poterit esse $+aRp$ sive $-aNp$.} (Casus sextus supra).

Etiam huius casus demonstrationem, quum prorsus similis sit demonstrationi casus tertii, brevitatis gratia omittimus.

\SourceLeaf{114}
\Article{141}

\textit{Casus quintus. Quando $T+1$ est formae $4n+3$ $(=b)$, $p$ eiusdem formae, atque $+pRb$ sive $-pNb$, nequit esse $+bRp$ sive $-bNp$.} (Casus tertius supra).

Sit $p\equiv e^2\ (\mathrm{mod.}\ b)$, atque $e$ par et $<b$.

I. Quando $e$ per $p$ non est divisibilis. Ponatur $e^2=p+bf$, eritque $f$ positivus, formae $4n+3$, $<b$ atque ad $p$ primus. Porro erit $pRf$ adeoque per prop. 13 art. 132, $-fRp$. Hinc et ex $+bfRp$ fit $-bRp$ adeoque $+bNp$. \textit{Q.~E.~D.}

II. Quando $e$ per $p$ est divisibilis, sit $e=pg$, atque $ggp=1+bh$. Tum erit $h$ formae $4n+1$ atque ad $p$ primus, $p\equiv g^2p^2\ (\mathrm{mod.}\ h)$, adeoque $pRh$; hinc fit $+hRp$ (prop. 10 art. 132), unde et ex $-bhRp$ sequitur $-bRp$, sive $+bNp$. \textit{Q.~E.~D.}

\Article{142}

\textit{Casus sextus. Quando $T+1$ est formae $4n+3$ $(=b)$, $p$ formae $4n+1$, atque $pRb$, non poterit esse $\pm bNp$.} (Supra casus septimus).

Demonstrationem praecedenti omnino similem omittimus.

\Article{143}

\textit{Casus septimus. Quando $T+1$ est formae $4n+3$ $(=b)$, $p$ eiusdem formae, atque $+pNb$ sive $-pRb$, non poterit esse $+bNp$ sive $-bRp$.} (Casus quartus supra).

Sit $-p\equiv e^2\ (\mathrm{mod.}\ b)$, atque $e$ par et $<b$.

I. Quando $e$ per $p$ non divisibilis. Sit $-p=e^2-bf$ eritque $f$ positivus, formae $4n+1$, ad $p$ primus ipsoque $b$ minor (etenim $e$ certo non maior quam $b-1$, $p<b-1$, quare erit $bf=e^2+p<b^2-b$, \textit{i. e.} $f<b-1$). Porro erit $-pRf$, hinc (prop. 10 art. 132) $+fRp$, unde et ex $+bfRp$ fit $+bRp$, sive $-bNp$.

II. Quando $e$ per $p$ est divisibilis, sit $e=pg$, atque $g^2p=-1+bh$. Tum erit $h$ positivus, formae $4n+3$, ad $p$ primus et $<b$. Porro erit $-pRh$, unde fit (prop. 14 art. 132) $+hRp$. Hinc et ex $bhRp$ sequitur $+bRp$ sive $-bNp$. \textit{Q.~E.~D.}

\SourceLeaf{115}
\Article{144}

\textit{Casus octavus. Quando $T+1$ est formae $4n+3$ $(=b)$, $p$ formae $4n+1$, atque $+pNb$ sive $-pRb$, non poterit esse $\pm bRp$.} (Casus ultimus supra).

Demonstratio perinde procedit ut in casu praecedenti.

\begin{center}
\textit{Methodus analoga, theorema art. 114 demonstrandi.}
\end{center}

\Article{145}

In demonstratt. praecc. semper pro $e$ valorem parem accepimus (artt. 137\,..\,144); observare convenit, etiam valorem imparem adhiberi potuisse, sed tum plures adhuc distinctiones introducendae fuissent. Qui his disquisitionibus delectantur, haud inutile facient, si vires suas in evolutione horum casuum exercitent. Praeterea theoremata ad residua $+2$ et $-2$ pertinentia tunc supponi debuissent; quum vero nostra demonstratio absque his theorematibus sit perfecta, novam hinc methodum nanciscimur, illa demonstrandi. Quae minime est contemnenda, quum methodi, quibus supra pro demonstratione theorematis, $\pm2$ esse residuum cuiusvis numeri primi formae $8n+1$, usi sumus, minus directae videri possint. Reliquos casus (qui ad numeros primos formarum $8n+3$, $8n+5$, $8n+7$ spectant) per methodos supra traditas demonstratos, illudque theorema tantummodo per inductionem inventum esse supponemus; hanc autem inductionem per sequentes reflexiones ad certitudinis gradum evehemus.

Si $\pm2$ omnium numerorum primorum formae $8n+1$ residuum non esset, ponatur minimus primus huius formae, cuius non-residuum $\pm2$, $=a$, ita ut pro omnibus primis ipso $a$ minoribus theorema valeat. Tum accipiatur numerus aliquis primus $<\frac{1}{2}a$, cuius non-residuum $a$ (qualem dari ex art. 129 facile deducitur). Sit hic $=p$ eritque per theor. fund. $pNa$. Hinc fit $\pm2pRa$. --- Sit itaque $e^2\equiv2p\ (\mathrm{mod.}\ a)$ ita ut $e$ sit impar atque $<a$. Tum duo casus erunt distinguendi.

I. Quando $e$ per $p$ non est divisibilis. Sit $e^2=2p+aq$ eritque $q$ positivus, formae $8n+7$ vel formae $8n+3$ (prout $p$ est formae $4n+1$ vel $4n+3$), $<a$, atque per $p$ non divisibilis. Iam omnes factores primi ipsius $q$ in quatuor classes distribuantur, sint scilicet $e$ formae $8n+1$, $f$ formae $8n+3$, $g$ formae $8n+5$, $h$ formae $8n+7$; productum e factoribus primae classis sit $E$, producta e factoribus secundae, tertiae, quartae classis respective, $F,G.H$\HistoricalNote{*)}{Si ex aliqua classe nulli factores adessent, loco producti ex his $1$ scribere oporteret.}.

\SourceLeaf{116}
His ita factis, consideremus \textit{primo} casum ubi $p$ est formae $4n+1$, sive $q$ formae $8n+7$. Tum facile perspicitur fore $2RE$, $2RH$, unde $pRE$, $pRH$, hincque tandem $ERp$, $HRp$. Porro erit $2$ non-residuum cuiusvis factoris formae $8n+3$ aut $8n+5$, adeoque etiam $p$; hinc quivis talis factor non-residuum ipsius $p$; unde facile concluditur $FG$ fore ipsius $p$ residuum si $f+g$ fuerit par, non-residuum si $f+g$ fuerit impar. At $f+g$ impar esse non potest; facile enim perspicietur omnes casus enumerando, $EFGH$ sive $q$ fieri vel formae $8n+3$ vel $8n+5$, si fuerit $f+g$ impar, quidquid sint singuli $e,f,g,h$, contra hyp. Erit igitur $FGRp$, $EFGHRp$, sive $qRp$, hincque tandem, propter $aqRp$, $aRp$ contra hyp. \textit{Secundo} quando $p$ est formae $4n+3$, simili modo ostendi potest, fore $pRE$ adeoque $ERp$, $-pRF$ adeoque $FRp$, tandem $g+h$ parem hincque $GHRp$, unde tandem sequitur $qRp$, $aRp$ contra hyp.

II. Quando $e$ per $p$ divisibilis, demonstratio simili modo adornari, et a peritis (quibus solis hic articulus est scriptus) haud difficulter evolvi poterit. Nos brevitatis gratia eam omittimus.

\begin{center}
\textit{Solutio problematis generalis.}
\end{center}

\Article{146}

Per theorema fundamentale atque propositiones ad residua $-1$ et $\pm2$ pertinentes semper determinari potest utrum numerus quicunque datus numeri primi dati residuum sit an non-residuum. At haud inutile erit, reliqua etiam quae supra tradidimus hic iterum in conspectum producere, ut omnia coniuncta habeantur quae sunt necessaria ad solutionem

\textsc{Problematis:} \textit{Propositis duobus numeris quibuscunque $P$, $Q$, invenire, utrum alter $Q$, alterius $P$ residuum sit an non-residuum.}

\textit{Sol.} I. Sit $P=a^{\alpha}b^{\beta}c^{\gamma}$ etc. designantibus $a,b,c$ etc. numeros primos inaequales positive acceptos (nam $P$ manifesto absolute est sumendus). Brevitatis gratia in hoc art. \textit{relationem} duorum numerorum $x,y$ simpliciter dicemus eam quatenus prior $x$ posterioris $y$ residuum est vel non-residuum. Pendet igitur relatio ipsorum $Q,P$ a relationibus ipsorum $Q,a^{\alpha}$; $Q,b^{\beta}$ etc. (art. 105).

II. Ut relatio ipsorum $Q,a^{\alpha}$ (de reliquis enim $Q,b^{\beta}$ etc. idem valet) innotescat, duo casus distinguendi.

\SourceLeaf{117}
1. Quando $Q$ per $a$ est divisibilis. Ponatur $Q=Q'a^e$, ita ut $Q'$ per $a$ non sit divisibilis. Tunc si $e=\alpha$ vel $e>\alpha$, erit $QRa^{\alpha}$; si vero $e<\alpha$ atque impar, erit $QNa^{\alpha}$: tandem si $e<\alpha$ atque par, habebit $Q$ ad $a^{\alpha}$ eandem relationem quam habet $Q'$ ad $a^{\alpha-e}$. Reductus est itaque hic casus ad

2. Quando $Q$ per $a$ non est divisibilis. Hic denuo duos casus distinguimus.

$(A)$ Quando $a=2$. Tunc semper erit $QRa^{\alpha}$, quando $\alpha=1$; quando vero $\alpha=2$, requiritur, ut sit $Q$ formae $4n+1$: denique quando $\alpha=3$ vel $>3$, $Q$ debet esse formae $8n+1$. Quae conditio si locum habet, erit $QRa^{\alpha}$.

$(B)$ Quando $a$ est alius numerus primus. Tunc $Q$ ad $a^{\alpha}$ eandem relationem habebit quam habet ad $a$ (V. art. 101).

III. Relatio numeri cuiuscunque $Q$ ad numerum primum $a$ (imparem) ita investigatur. Quando $Q>a$, substituatur loco ipsius $Q$ ipsius residuum minimum positivum secundum modulum $a$\HistoricalNote{*)}{\textit{Residuum} in signific. art. 4. --- Plerumque praestat residuum \textit{absolute} minimum accipere.}. Hoc ad $a$ eandem relationem habebit quam habet $Q$.

Porro resolvatur $Q$, sive numerus ipsius loco assumtus, in factores suos primos $p,p',p''$ etc., quibus adiungendus factor $-1$, quando $Q$ est negativus. Tum constat relationem ipsius $Q$ ad $a$ pendere a relationibus singulorum $p,p',p''$ etc. ad $a$. Scilicet si inter illos factores sunt $2m$ non-residua ipsius $a$, erit $QRa$, si vera $2m+1$, erit $QNa$. Facile autem perspicitur, si inter factores $p,p',p''$ etc., bini aut quaterni aut seni aut generaliter $2k$ aequales occurrant, hos tuto eiici posse.

IV. Si inter factores $p.p',p''$ reperiuntur $-1$ et $2$, horum relatio ad $a$ ex artt. 108, 112, 113, 114 inveniri potest. Reliquorum autem relatio ad $a$ pendet a relatione ipsius $a$ ad ipsos (theor. fund., atque propp. art. 131). Sit $p$ unus ex ipsis, invenieturque, (tractando numeros $a,p$ eodem modo ut antea $Q$ et $a$ illis respective maiores) relationem ipsius $a$ ad $p$ aut per artt. 108---114 determinari posse (si scilicet residuum minimum ipsius $a$ $(\mathrm{mod.}\ p)$ nullos factores primos impares habeat), aut insuper a relatione ipsius $p$ ad numeros quosdam primos ipso $p$ minores pendere. Idem valet de reliquis factoribus $p',p''$ etc. Facile iam \SourceLeaf{118}perspicitur per continuationem huius operationis tandem ad numeros perventum iri quorum relationes per propp. artt. 108---114 determinari possint. Per exemplum haec clariora fient.

\textit{Ex.} Quaeritur relatio numeri $+453$ ad $1236$. Est $1236=4.3.103$; $+453R4$ per II. 2 $(A)$; $+453R3$ per II. 1. Superest igitur ut relatio ipsius $+453$ ad $103$ exploretur. Eadem autem erit quam habet $+41$ $(\equiv453,\ \mathrm{mod.}\ 103)$ ad $103$; eadem ipsius $+103$ ad $41$ (theor. fund.), sive ipsius $-20$ ad $41$. At est $-20R41$; namque $-20=-1.2.2.5$; $-1R41$ (art. 108); atque $+5R41$ ideo quod $41\equiv1$ adeoque ipsius $5$ residuum est (theor. fund.). Hinc sequitur $+453R103$, hincque tandem $+453R1236$. Est autem revera $453\equiv297^2\ (\mathrm{mod.}\ 1236)$.

\begin{center}
\textit{De formis linearibus omnes numeros primos continentibus, quorum vel residuum vel non-residuum est numerus quicunque datus.}
\end{center}

\Article{147}

Proposito numero quocunque $A$, \textit{formulae} certae exhiberi possunt, sub quibus omnes numeri ad $A$ primi quorum residuum est $A$ continentur, sive omnes qui esse possunt \textit{divisores} numerorum formae $xx-A$ (designante $xx$ quadratum indeterminatum)\HistoricalNote{*)}{Huiusmodi numeros simpliciter \textit{divisores ipsius $xx-A$} dicemus unde sponte patet quid sint \textit{non-divisores}.}. Sed brevitatis gratia ad eos tantum divisores respiciemus, qui sunt impares atque ad $A$ primi, quum ad hos casus reliqui facile reduci possint.

Sit primo $A$ aut numerus primus positivus formae $4n+1$, aut negativus formae $4n-1$. Tum secundum theorema fundamentale omnes numeri primi, qui, positive sumti, sunt residua ipsius $A$, erunt divisores ipsius $xx-A$: omnes autem numeri primi (excepto numero $2$ qui semper est divisor) qui ipsius $A$ sunt non-residua erunt non-divisores ipsius $xx-A$. Sint omnia residua ipsius $A$ ipso $A$ minora (exclusa cifra) $r,r',r''$ etc. omnia non-residua vero $n,n',n''$ etc. Tum quivis numerus primus, in aliqua formarum $Ak+r$, $Ak+r'$, $Ak+r''$ etc. contentus, erit divisor ipsius $xx-A$, quivis autem primus in aliqua formarum $Ak+n$, $Ak+n'$ etc. contentus non-divisor erit, designante $k$ numerum integrum indeterminatum. Illas formas dicimus \textit{formas divisorum ipsius $xx-A$}, has vero \textit{formas non-divisorum}. Utrorumque multitudo erit $\frac{1}{2}(A-1)$. Porro si $B$ est numerus compositus impar atque $ARB$, omnes factores primi ipsius $B$ in aliqua for\SourceLeaf{119}marum priorum continentur adeoque etiam $B$. Quare \textit{quivis} numerus impar in forma non-divisorum contentus, erit non-divisor formae $xx-A$. Sed hoc theorema convertere non licet; nam si $B$ est non-divisor compositus impar formae $xx-A$, inter factores primos ipsius $B$ aliqui non-divisores erunt, quorum multitudo si est \textit{par}, $B$ nihilominus in aliqua forma divisorum reperietur. V. art. 99.

\textit{Ex.} Hoc modo pro $A=-11$ formae divisorum ipsius $xx+11$ inveniuntur hae: $11k+\nobreak{}1,\allowbreak{}3,\allowbreak{}4,\allowbreak{}5,\allowbreak{}9$; formae non-divisorum autem erunt $11k+\nobreak{}2,\allowbreak{}6,\allowbreak{}7,\allowbreak{}8,\allowbreak{}10$. Erit itaque $-11$ non-residuum omnium numerorum imparium, qui in aliqua posteriorum formarum continentur, residuum autem omnium primorum ad aliquam priorum pertinentium.

Similes formae dantur pro divisoribus atque non-divisoribus ipsius $xx-A$, quemcunque numerum designet $A$. Sed facile perspicitur, eos ipsius $A$ valores tantummodo considerari oportere, qui per nullum quadratum sint divisibiles; patet enim si fuerit $A=a^2A'$, omnes divisores\HistoricalNote{*)}{Nempe qui sint primi ad $A$.} ipsius $xx-A$ etiam fore divisores ipsius $xx-A'$, similiterque non-divisores. --- Distinguemus autem tres casus, 1) quando $A$ est formae $+(4n+1)$ vel $-(4n-1)$. 2) quando $A$ est formae $-(4n+1)$ vel $+(4n-1)$. 3) quando $A$ est par sive formae $\pm(4n+2)$.

\Article{148}

\textit{Casus primus}, quando $A$ est formae $+(4n+1)$ vel $-(4n-1)$. Resolvatur $A$ in factores suos primos, tribuaturque iis qui sunt formae $4n+1$ signum positivum, iis vero qui sunt formae $4n-1$ signum negativum (unde fiet productum ex ipsis $=A$). Sint hi factores $a,b,c,d$ etc. Distribuantur omnes numeri ipso $A$ minores et ad $A$ primi in duas classes, et quidem in primam classem omnes numeri qui sunt nullius ex numeris $a,b,c,d$ etc. non-residua, aut duorum, aut quatuor aut generaliter multitudinis paris; in secundam vero ii, qui sunt non-residua unius ex numeris $a,b,c$ etc. aut trium etc. aut generaliter multitudinis imparis. Designentur priores per $r,r',r''$ etc., posteriores per $n,n',n''$ etc. Tum formae $Ak+r$, $Ak+r'$, $Ak+r''$ etc. erunt formae divisorum ipsius $xx-A$, formae vero $Ak+n$, $Ak+n'$ etc. erunt formae non-divisorum ipsius $xx-A$ (\textit{i. e. numerus quicunque primus, praeter $2$, erit divisor aut non-divisor ipsius $xx-A$, prout in aliqua formarum priorum aut posteriorum continetur}). Si enim $p$ est nume\SourceLeaf{120}rus primus positivus atque alicuius ex numeris $a,b,c$ etc. residuum vel non-residuum, hic ipse numerus ipsius $p$ residuum vel non-residuum erit (theor. fund.). Quare si inter numeros $a,b,c$ etc. sunt $m$, quorum non-residuum est $p$, totidem erunt non-residua ipsius $p$, adeoque si $p$ in aliqua formarum priorum continetur, erit $m$ par et $ARp$, si vero in aliqua posteriorum, erit $m$ impar atque $ANp$.

\textit{Ex.} Sit $A=+105=-3\times+5\times-7$. Tum numeri $r,r',r''$ etc. erunt hi: $1,\allowbreak{}4,\allowbreak{}16,\allowbreak{}46,\allowbreak{}64,\allowbreak{}79$ (qui sunt non-residua nullius numerorum $3,\allowbreak{}5,\allowbreak{}7$); $2,\allowbreak{}8,\allowbreak{}23,\allowbreak{}32,\allowbreak{}53,\allowbreak{}92$ (qui sunt non-residua numerorum $3,\allowbreak{}5$); $26,\allowbreak{}41,\allowbreak{}59,\allowbreak{}89,\allowbreak{}101,\allowbreak{}104$ (qui sunt non-residua numerorum $3,\allowbreak{}7$); $13,\allowbreak{}52,\allowbreak{}73,\allowbreak{}82,\allowbreak{}97,\allowbreak{}103$ (qui sunt non-residua numerorum $5,\allowbreak{}7$). --- Numeri autem $n,n',n''$ etc. erunt hi: $11,\allowbreak{}29,\allowbreak{}44,\allowbreak{}71,\allowbreak{}74,\allowbreak{}86$; $22,\allowbreak{}37,\allowbreak{}43,\allowbreak{}58,\allowbreak{}67,\allowbreak{}88$; $19,\allowbreak{}31,\allowbreak{}34,\allowbreak{}61,\allowbreak{}76,\allowbreak{}94$; $17,\allowbreak{}38.47,\allowbreak{}62,\allowbreak{}68,\allowbreak{}83$. Seni primi sunt non-residua ipsius $3$, seni posteriores non-residua ipsius $5$, tum sequuntur non-residua ipsius $7$, tandem ii qui sunt non-residua omnium trium simul.

Facile ex combinationum theoria atque artt. 32, 96 deducitur, numerorum $r,r',r''$ etc. multitudinem fore
\[
=t\left(1+\frac{l\, .\, l-1}{1\, .\, 2}+\frac{l\, .\, l-1\, .\, l-2\, .\, l-3}{1\, .\, 2\, .\, 3\, .\, 4}+ .\, .\, .\, .\, .\right)
\]
numerorum $n,n',n''$ etc. multitudinem
\[
=t\left(l+\frac{l\, .\, l-1\, .\, l-2}{1\, .\, 2\, .\, 3}+\frac{l\, .\, l-1\, .\, .\, .\, .\, l-4}{1\, .\, 2\, .\, .\, .\, .\, 5}+ .\, .\, .\right)
\]
ubi $l$ designat multitudinem numerorum $a,b,c$ etc.;
\[
t=2^{-l}(a-1)(b-1)(c-1)\text{ etc.}
\]
et utraque series continuanda donec abrumpatur. (Dabuntur scilicet $t$ numeri qui sunt residua omnium $a,b,c$ etc., $\frac{t\, .\, l\, .\, l-1}{1\, .\, 2}$ qui sunt non-residua duorum, etc. sed demonstrationem hanc fusius explicare brevitas non permittit). Utriusque autem seriei summa\HistoricalNote{*)}{Neglecto factor\'e $t$.} est $=2^{l-1}$. Scilicet prior prodit ex hac
\[
1+(l-1)+\frac{l-1\, .\, l-2}{1\, .\, 2}+ .\, .\, .
\]
iungendo terminum secundum et tertium, quartum et quintum etc., posterior vero ex eadem iungendo terminum primum atque secundum, tertium et quartum etc. Dabuntur itaque tot formae divisorum ipsius $xx-A$, quot dantur formae non-divisorum, scilicet $\frac{1}{2}(a-1)(b-1)(c-1)$ etc.

\SourceLeaf{121}
\Article{149}

\textit{Casum secundum et tertium} hic simul contemplari possumus. Poterit scilicet $A$ semper hic poni $=(-1)Q$, aut $=(+2)Q$, aut $=(-2)Q$, designante $Q$ numerum formae $+(4n+1)$, aut $-(4n-1)$, quales in art. praec. consideravimus. Sit generaliter $A=\alpha Q$, ita ut sit $\alpha$ aut $=-1$, aut $=\pm2$. Tum erit $A$ residuum omnium numerorum, quorum residuum est aut uterque $\alpha$ et $Q$, aut neuter; non-residuum autem omnium, quorum non-residuum alteruter tantum numerorum $\alpha,Q$. Hinc formae divisorum ac non-divisorum ipsius $xx-A$ facile derivantur. Si $\alpha=-1$, distribuantur omnes numeri ipso $4A$ minores ad ipsumque primi in duas classes, in priorem ii, qui sunt in aliqua forma divisorum ipsius $xx-Q$ simulque in forma $4n+1$, iique, qui sunt in aliqua forma non-divisorum ipsius $xx-Q$ simulque in forma $4n+3$; in posteriorem reliqui. Sint priores $r,r',r''$ etc., posteriores $n,n',n''$ etc., eritque $A$ residuum omnium numerorum primorum in aliqua formarum $4Ak+r$, $4Ak+r'$, $4Ak+r''$ etc. contentorum, non-residuum autem omnium primorum in aliqua formarum $4Ak+n$, $4Ak+n'$ etc. contentorum. --- Si $\alpha=\pm2$, distribuantur omnes numeri ipso $8Q$ minores ad ipsumque primi in duas classes, in primam ii, qui continentur in aliqua forma divisorum ipsius $xx-Q$ simulque in aliqua formarum $8n+1$, $8n+7$ pro signo superiori, vel formarum $8n+1$, $8n+3$ pro inferiori, iique qui contenti sunt in aliqua forma non-divisorum ipsius $xx-Q$ simulque in aliqua harum $8n+3$, $8n+5$ pro signo superiori, vel harum $8n+5$, $8n+7$ pro inferiori, --- in secundam reliqui. Tum designatis numeris classis prioris per $r,r',r''$ etc., numerisque classis posterioris per $n,n',n''$ etc., $\pm2Q$ erit residuum omnium numerorum primorum in aliqua formarum $8Qk+r$, $8Qk+r'$, $8Qk+r''$ etc. contentorum, omnium autem primorum in aliqua formarum $8Qk+n$, $8Qk+n'$, $8Qk+n''$ etc. non-residuum. Ceterum facile demonstrari potest, etiam hic totidem formas divisorum ipsius $xx-A$ datum iri ac non-divisorum.

\textit{Ex.} Hoc modo invenitur $+10$ esse residuum omnium numerorum primorum in aliqua formarum $40k+\nobreak{}1,\allowbreak{}3,\allowbreak{}9,\allowbreak{}13,\allowbreak{}27,\allowbreak{}31,\allowbreak{}37,\allowbreak{}39$ contentorum, non-residuum vero omnium primorum, qui sub aliqua formarum $40k+\nobreak{}7,\allowbreak{}11,\allowbreak{}17,\allowbreak{}19,\allowbreak{}21,\allowbreak{}23,\allowbreak{}29,\allowbreak{}33$ continentur.
\par\endgroup
